\documentclass[11pt,a4paper]{article}

\usepackage[utf8]{inputenc}
\usepackage[T1]{fontenc}
\usepackage{lmodern}
\usepackage{microtype}
\usepackage{amsmath,amssymb,amsthm}
\usepackage{mathtools}
\usepackage{geometry}
\usepackage{hyperref}

\geometry{a4paper,margin=25mm}

\newtheorem{theorem}{Theorem}[section]

\title{\textbf{Identity Bedrock: A Minimal Cryptographic Substrate and Its Systemic Implications}}
\author{Jean-François Bouchard}
\date{March 2026}

\begin{document}

\maketitle

\begin{abstract}
We present Identity Bedrock, a deployed identity system constructed from a minimal cryptographic substrate. Starting from the requirement that no institutional key exists beneath the individual, we derive an architecture grounded in one-way functions (OWFs), which are necessary and sufficient for existentially unforgeable signature schemes under adaptive chosen-message attack.

The system separates identity cryptography from directory consensus via a chain-agnostic, append-only binding layer, enabling portability across heterogeneous substrates while preserving verifiability.

The system has been deployed in production for over five years, with publicly verifiable records and continuous operation.

We show that additional system capabilities—including naming, structured representation, content retrieval, execution, connectivity, and economic interaction—emerge as derived properties of this identity substrate rather than independent subsystems.
\end{abstract}

\section{Introduction}

Digital identity systems commonly rely on institutional trust anchors such as certificate authorities, identity providers, or registry operators. These introduce external control points over identity and create dependencies that affect availability, portability, and sovereignty.

This work describes a deployed system rather than a purely theoretical construction. It explores how identity can be constructed by minimizing such dependencies and relying only on necessary cryptographic assumptions.

\subsection{Contributions}

\begin{itemize}
\item Minimal construction grounded in OWF-based signatures
\item Chain-agnostic identity and directory separation
\item Publicly deployed and verifiable system
\item Demonstration that system-level capabilities emerge from identity
\end{itemize}

\section{Background}

One-way functions (OWFs) are necessary and sufficient for existentially unforgeable signature schemes under adaptive chosen-message attack. This provides a minimal foundation for authentication.

Traditional identity systems introduce hierarchical trust models that extend beyond this minimal requirement, adding operational dependencies.

\section{Architecture}

\subsection{Identity Cryptography}

A user generates a keypair $(sk, pk)$ from a high-entropy seed. Authentication consists of signing a challenge and verifying the signature.

\subsection{Directory Substrate}

An append-only directory binds identifiers to public keys, providing:

\begin{itemize}
\item Global verifiability
\item Persistence
\item Finality
\end{itemize}

Identity and directory layers are decoupled, allowing deployment across different substrates.

\section{Security}

\begin{theorem}
If the signature scheme is existentially unforgeable under adaptive chosen-message attack and the directory is append-only with finality, then authentication is unforgeable under standard assumptions.
\end{theorem}

\begin{proof}
A successful forgery implies a break of the signature scheme or a violation of directory integrity.
\end{proof}

\section{Derived System Properties}

\subsection{Naming and Resolution}

Identity is bound to naming through an append-only name-value substrate (EmerDNS), where records are signed by the controlling key:

\begin{equation}
(n, v, pk, \sigma), \quad \sigma = \text{Sign}_{sk}(n \parallel v)
\end{equation}

Resolution consists of retrieving $(v, pk, \sigma)$ and verifying $\sigma$ against $pk$.

The system integrates with standard DNS resolution paths and adheres to RFC~1035 semantics. Resolution may be performed through DNS servers implementing EmerDNS or via operating system-level name resolution policies.

\subsection{Structured Representation}

Objects are represented using deterministic formats, enabling consistent parsing, indexing, and decentralized query.

\subsection{Content Addressing}

Content is retrieved via cryptographic hashes, ensuring integrity independent of transport.

\subsection{Deterministic Execution}

\begin{equation}
Execution = Hash(Input, Environment, Command)
\end{equation}

This enables reproducible and verifiable execution.

\subsection{Connectivity}

Overlay networking enables communication independent of underlying network constraints.

\subsection{Economic Interaction}

Incentive mechanisms operate through identity-bound attribution without centralized coordination.

\section{Deployment and Operation}

The system is deployed and operational. Identity bindings and naming records are actively in use and publicly verifiable.

\subsection{Public Verification}

A representative deployment is accessible at:

\begin{center}
\url{https://explorer.emercoin.com/dpo/PrivateNESS.Network}
\end{center}

Any participant may independently retrieve and verify records.

\subsection{Application Resolution}

Applications are resolved through identity-bound naming, retrieved via content-addressing, and executed deterministically.

\subsection{Node Deployment}

Nodes operate on commodity hardware, including low-cost and repurposed devices, without requiring specialized infrastructure.

\section{Implementation}

The system is implemented using Emercoin's Name-Value Storage (EmerDNS), anchored via AuxPoW.

Identity bindings are stored as signed records and are publicly verifiable through standard node software.

Additional components operate on Fiber-based networks using compatible key derivation and signature mechanisms.

The system integrates with existing protocols, including DNS, SSH, and TLS workflows, without requiring modification to client software.

\section{Evaluation}

Authentication requires only signature verification and directory lookup, resulting in low computational overhead.

The system has operated continuously in production environments, demonstrating stability over multiple years.

\section{Discussion}

The architecture demonstrates that naming, execution, connectivity, and economic behavior can be understood as consequences of a unified identity substrate.

\section{Conclusion}

Identity Bedrock demonstrates that a minimal cryptographic substrate combined with an append-only directory is sufficient to support both authentication and broader system functionality.

\end{document}